\par\medskip
\begin{center}\hypertarget{native-164-article117}{}\pdfbookmark[2]{117.}{article-117}117.\end{center}
The position of the Earth computed in this way from the tables must be referred to the center of the Earth, whereas the observed position of the celestial body is referred to a point on the Earth's surface. There are three ways to remedy this discrepancy. The observation can be reduced to the center of the Earth, that is, corrected for parallax. Or the heliocentric position of the Earth can be reduced to the actual place of observation by properly applying parallax to the position of the Sun computed from the tables. Finally, both positions can be transferred to some third point, which is most conveniently taken to be the intersection of the line of sight with the plane of the ecliptic. The observation itself is then retained as made. We showed how to reduce the position of the Earth to this point in Article 72. The first method can be used only if the distance of the celestial body from the Earth is known at least approximately. It is then quite convenient, especially when the observation is made on the meridian itself, where only the declination is affected by parallax. It is preferable, however, to apply this reduction directly to the observed position before undertaking the transformations in the preceding article. If the distance from the Earth is still entirely unknown, we must resort to the second or third method. The second is used whenever the equator is taken as the fundamental plane, and the third is preferred whenever all positions are to be referred to the ecliptic.
